\documentclass{article}
\usepackage[T1]{fontenc}
\usepackage{lmodern}
\usepackage{graphicx}
\usepackage{amsmath,amssymb}
\usepackage[a4paper,margin=2cm]{geometry}
\usepackage[absolute,overlay]{textpos}
\setlength{\TPHorizModule}{1mm}
\setlength{\TPVertModule}{1mm}
\usepackage[indonesian]{babel}
\usepackage{hyperref}
\setlength{\emergencystretch}{2em}
% unit: Halaman 16

\begin{document}

% === Halaman 16 (python/native) ===

16

4. Adaptive-chosen-plaintext attack

• Adaptive chosen plaintext attack adalah variasi dari chosen plaintext attack, 

dalam hal ini penyerang tidak hanya memilih plainteks, tapi dapat 

menyesuaikan plainteks berikutnya berdasarkan hasil ciphertext dari query 

sebelumnya.

• Artinya, penyerang belajar secara bertahap: setelah melihat hasil enkripsi 

pertama, ia memilih plainteks baru yang lebih “cerdas” untuk menggali informasi 

lebih dalam tentang algoritma/kunci.

\end{document}
